\subsection{伴随素理想的局部化}

命题 5. 设 A 是环，S 是 A 的乘性子集，\(\Phi\) 是 A 中不与 S 相交的素理想集合，M 是 A-模。则：

\begin{enumerate}
\def\labelenumi{(\roman{enumi})}
\item
  \(\mathfrak{p} \mapsto \mathrm{S}^{-1}\mathfrak{p}\) 是从 \(\operatorname{Ass}_{\mathbf{A}}(\mathbf{M}) \cap \Phi\) 到 \(\operatorname{Ass}_{\mathbf{S}^{-1}\mathbf{A}}(\mathbf{S}^{-1}\mathbf{M})\) 的某个子集的双射。
\item
  如果 \(\mathfrak{p} \in \Phi\) 是有限生成理想，且 \(S^{-1}\mathfrak{p} \in \operatorname{Ass}_{\mathbf{S}^{-1}\mathbf{A}}(\mathbf{S}^{-1}\mathbf{M})\)，则 \(\mathfrak{p} \in \operatorname{Ass}_A(M)\)。
\end{enumerate}

回忆（第二章，第 2 节，第 5 节，命题 11），映射\(\mathfrak{p} \mapsto \mathrm{S}^{-1}\mathfrak{p}\)是从\(\Phi\)到\(\mathrm{S}^{-1}\mathrm{A}\)的素理想集合的双射。如果\(\mathfrak{p} \in \operatorname{Ass}_{\mathbf{A}}(\mathbf{M}) \cap \Phi\)，则\(\mathfrak{p}\)是\(\mathbf{N}\)的一个单生成子模\(\mathbf{M}\)的湮灭子；于是\(\mathrm{S}^{-1}\mathfrak{p}\)是\(\mathrm{S}^{-1}\mathrm{N}\)的单生成子模\(\mathrm{S}^{-1}\mathrm{M}\)的湮灭子（第二章，第 2 节，第 4 节，公式 (9)），从而\(\mathbf{S}^{-1}\mathfrak{p}\in\mathrm{Ass}_{\mathbf{S}^{-1}\mathbf{A}}(\mathbf{S}^{-1}\mathbf{M})\)。反之，设\(\mathfrak{p}\in\Phi\)是有限生成的，且\(\mathbf{S}^{-1}\mathfrak{p}\)与\(\mathbf{S}^{-1}\mathbf{M}\)相伴；则存在\(x\in\mathbf{M}\)和\(t\in\mathbf{S}\)，使得\(\mathbf{S}^{-1}\mathfrak{p}\)是\(x/t\)的湮灭子。设\((a_{i})_{1\leqslant i\leqslant n}\)是\(\mathfrak{p}\)的一组生成元；则\((a_{i}/1)(x/t)=0\)，从而存在\(s_{i}\in\mathbf{S}\)，使得\(s_{i}a_{i}x=0\)（\(1\leqslant i\leqslant n\)）。记\(s=s_{1}s_{2}\ldots s_{n}\)；对于所有\(a\in\mathfrak{p}\)，有\(sax=0\)，从而\(\mathfrak{p}\subset\mathrm{Ann}(sx)\)；另一方面，若\(b\in\mathbf{A}\)满足\(bsx=0\)，则按定义\(b/1\in\mathbf{S}^{-1}\mathfrak{p}\)，从而\(b\in\mathfrak{p}\)。于是\(\mathfrak{p}=\mathrm{Ann}(sx)\)，且\(\mathfrak{p}\in\mathrm{Ass}_{\mathbf{A}}(\mathbf{M})\)。

推论。若环 A 是诺特环，则映射 \(\mathfrak{p} \mapsto \mathrm{S}^{-1}\mathfrak{p}\) 是从 \(\operatorname{Ass}_{\mathbf{A}}(\mathbf{M}) \cap \Phi\) 到 \(\operatorname{Ass}_{\mathbf{S}^{-1}\mathbf{A}}(\mathbf{S}^{-1}\mathbf{M})\) 的双射。

若 A 不是诺特环，则从 \(p \mapsto S^{-1}p\) 的 \(\operatorname{Ass}_{\mathbf{A}}(\mathbf{M}) \cap \Phi\) 到 \(\operatorname{Ass}_{\mathbf{S}^{-1}\mathbf{A}}(\mathbf{S}^{-1}\mathbf{M})\) 的映射不一定是满射（习题 1）。

命题 6. 设 A 是诺特环，M 是 A-模，S 是 A 的乘性子集，\(\Psi\) 是 \(\operatorname{Ass}_{\mathbf{A}}(\mathbf{M})\) 中不与 S 相交的元素所成的集合。则典范映射 \(M \to S^{-1}M\) 的核 N 是 M 的唯一子模，满足下列关系

\[
\operatorname{Ass} (N) = \operatorname{Ass} (M) - \Psi , \quad \operatorname{Ass} (M / N) = \Psi .\tag{3}
\]

根据第 1 节命题 4，存在\(\mathbf{N}'\)是\(\mathbf{M}\)的子模，满足关系\(\operatorname{Ass}(\mathbf{N}') = \operatorname{Ass}(\mathbf{M}) - \Psi'\)以及\(\operatorname{Ass}(\mathbf{M}/\mathbf{N}') = \Psi'\)。我们需要证明\(\mathbf{N}' = \mathbf{N}\)。考虑交换图

\[
\begin{array}{c} \mathrm{M} \xrightarrow {p} \mathrm{M} / \mathrm{N} ^ {\prime} \\ u \Biggl \downarrow \\ \mathrm{S} ^ {- 1} \mathrm{M} \xrightarrow [ \mathrm{S} ^ {- 1 p} ]{} \mathrm{S} ^ {- 1} (\mathrm{M} / \mathrm{N} ^ {\prime}) \end{array}
\]

其中 \(p, u, v\) 是典范同态。我们将证明 \(S^{-1}p\) 和 \(v\) 是单射；这将证明 \(u\) 和 \(p\) 具有相同的核，从而 \(N' = N\)。

由于\(\operatorname{Ass}(\mathbf{N}') \cap \Psi' = \varnothing\)，\(\operatorname{Ass}(\mathbf{N}')\)的每个元素都与 S 相交。于是\(\operatorname{Ass}_{\mathbf{S}^{-1}\mathbf{A}}(\mathbf{S}^{-1}\mathbf{N}') = \varnothing\)（命题 5 的推论），从而\(\mathbf{S}^{-1}\mathbf{N}' = \{0\}\)（第 1 节命题 2 的推论 1），这说明\(\mathbf{S}^{-1}\mathfrak{p}\)是单射（第二章，第 2 节，第 4 节，定理 1）。另一方面，若\(x\)属于核 K 的\(v\)，则\(\operatorname{Ann}(x) \cap \mathbf{S} \neq \varnothing\)（第二章，第 2 节，第 2 节，命题 4）；于是\(\operatorname{Ass}(\mathbf{K}) = \varnothing\)，因为\(\operatorname{Ass}(\mathbf{K}) \subset \operatorname{Ass}(\mathbf{M}/\mathbf{N}') = \Psi'\)；我们推出\(\mathbf{K} = \{0\}\)（第 1 节命题 2 的推论 1），且\(v\)是单射。

